
\section{LITERATURE SURVEY}

\hspace*{0.3in} The rapid proliferation of cloud-native architectures and microservice-based deployments has fundamentally transformed the operational landscape of modern software systems. Traditional monitoring approaches, rooted in threshold-based alerting and rule-driven correlation, were designed for monolithic or small-scale distributed environments where component interdependencies were well-understood and relatively static. As organizations began decomposing applications into hundreds or thousands of independently deployable services, these legacy strategies proved inadequate in diagnosing complex, cascading failures. Early AIOps frameworks, as surveyed by \cite{dang2019aiops}, attempted to bridge this gap by integrating statistical anomaly detection and log-based correlation into existing monitoring pipelines. However, such solutions remained reactive in nature, relying on manually curated rules and failing to generalize across evolving service topologies. The need for a fundamentally data-driven, topology-aware approach to failure diagnosis in microservice ecosystems became increasingly evident, motivating the design of more sophisticated AI-driven platforms capable of operating at scale under dynamic, production-grade conditions \cite{chen2020towards}.\\\\

\hspace*{0.3in} Graph Neural Networks (GNNs) have emerged as a compelling paradigm for modeling structured relational data, making them particularly suitable for capturing inter-service dependencies within microservice architectures. Zhao et al. \cite{zhao2021graph} demonstrated that service call graphs, when encoded as heterogeneous attributed graphs and processed through message-passing GNNs, enable anomaly detection that is sensitive to both local metric deviations and global topological context. Building on this, the use of Variational Autoencoders (VAEs) in conjunction with GNNs has enabled unsupervised anomaly detection by learning low-dimensional latent representations of normal service behavior; deviations from the learned distribution signal potential faults \cite{li2022vae}. Ma et al. \cite{ma2023gnn} further showed that temporal GNNs can detect transient anomalies in streaming metric data by modeling time-evolving service graphs. Despite these advances, most GNN-based detectors are designed as isolated components and do not feed seamlessly into downstream causal inference or remediation pipelines, a limitation that SYNAPSE explicitly addresses by embedding a VAE-GNN anomaly detector as the third stage of its integrated ten-stage AI pipeline.\\\\

\hspace*{0.3in} Establishing causal relationships among failing microservices — as opposed to merely identifying statistical correlations — is critical to accurate root cause analysis. Classical approaches such as the PC algorithm \cite{spirtes2000causation} and the Fast Causal Inference (FCI) algorithm operate on observational data to recover the skeleton of a causal Directed Acyclic Graph (DAG), but they are sensitive to distributional assumptions and scale poorly to high-dimensional service graphs with hundreds of concurrent metrics. Microsoft's CausaLM \cite{eberts2022causalm} and Alibaba's MicroCause \cite{wu2021microrca} represent industry efforts to adapt causal inference to microservice root cause localization, using Granger causality and propagation-based correlation respectively. More recently, the Differentiable Causal Discovery from Interventional data (DECI) framework \cite{geffner2022deep} introduced an end-to-end differentiable approach to learning causal graphs, treating causal structure learning as a continuous optimization problem. SYNAPSE adopts DECI as its causal discovery engine, augmenting it with a Retrieval-Augmented Generation (RAG)-constrained initialization prior ($W_0$) derived from historical incident patterns, thereby grounding learned causal structures in domain-specific operational knowledge rather than relying purely on observational signals.\\\\

\hspace*{0.3in} The application of Large Language Models (LLMs) to system observability and failure explanation constitutes an emerging research frontier. Ahmed et al. \cite{ahmed2023log} showed that pre-trained language models fine-tuned on structured log sequences can classify failure modes with high accuracy, while also generating natural-language summaries comprehensible to on-call engineers. Google's work on AI-assisted incident management \cite{google2023incident} demonstrated that LLMs can synthesize alert streams, runbooks, and historical postmortems to recommend remediation steps in real time. Root cause narratives generated by GPT-4-class models have been shown to reduce mean time to resolution (MTTR) by enabling junior engineers to act on AI-generated explanations without deep system expertise \cite{chen2023llm}. A key limitation in prior work, however, is that LLM outputs are not grounded in structured causal evidence — models may hallucinate plausible-sounding but factually incorrect diagnoses. SYNAPSE mitigates this by feeding structured outputs from the causal discovery and root cause extraction stages directly into the GPT-4o-mini reasoning prompt, ensuring that natural-language explanations are anchored to verified causal chains rather than generated purely from parametric model knowledge.\\\\

\hspace*{0.3in} Continual learning — the ability of ML models to incorporate new knowledge without catastrophically forgetting previously acquired representations — is indispensable in production AIOps systems, where service behavior, deployment configurations, and failure modes evolve continuously over time. Kirkpatrick et al. \cite{kirkpatrick2017ewc} introduced Elastic Weight Consolidation (EWC), a regularization-based approach that selectively protects weights critical to prior tasks during training on new data, effectively mitigating catastrophic forgetting in neural networks. Complementarily, Model-Agnostic Meta-Learning (MAML) \cite{finn2017maml} enables models to learn generalizable initializations from a distribution of tasks, allowing rapid adaptation to novel failure scenarios with minimal labeled examples — a property essential in low-signal environments following rare but critical system faults. Experience replay strategies, as applied in the continual learning literature \cite{rolnick2019experience}, maintain a memory buffer of representative past samples to prevent distributional drift. SYNAPSE's continual learning stage integrates EWC regularization with MAML-based meta-learning across both the VAE-GNN anomaly detector and the severity classifier, ensuring the platform remains accurate as microservice topologies evolve without requiring complete model retraining.\\\\

\hspace*{0.3in} Several commercial AIOps platforms have been developed to address observability and incident management in distributed systems, including Dynatrace, Datadog, New Relic, Moogsoft, and PagerDuty AIOps. Dynatrace's Davis AI engine employs a topology-aware anomaly detection mechanism based on Smartscape, its dynamic service dependency map, combined with deterministic causal inference to identify root causes \cite{dynatrace2022davis}. Datadog's Watchdog uses automated anomaly detection over metrics and APM traces, surfacing potential issues without manual threshold configuration \cite{datadog2023watchdog}. Moogsoft applies clustering and correlation algorithms to noise-reduce alert streams and identify incident patterns \cite{moogsoft2021aiops}. Despite their maturity, these platforms share critical limitations: they operate as closed, vendor-specific ecosystems with limited extensibility; their causal reasoning is largely heuristic rather than grounded in formal causal inference; they provide no explicit severity classification beyond basic priority scoring; and their remediation recommendations are static, drawn from predefined runbooks rather than dynamically generated and adapted. Crucially, none of these platforms offer continual learning capabilities that allow their AI models to evolve in response to new operational patterns without manual retraining cycles.\\\\

\hspace*{0.3in} Retrieval-Augmented Generation (RAG) \cite{lewis2020rag} has demonstrated substantial promise in enhancing the factual grounding and domain specificity of LLM outputs by conditioning generation on dynamically retrieved, contextually relevant documents. In the AIOps context, RAG architectures enable LLM-based reasoning engines to retrieve and incorporate historical incident reports, runbooks, architecture documentation, and past RCA outcomes at inference time, substantially reducing hallucination risk and improving recommendation quality. Jin et al. \cite{jin2023rag} applied RAG to log-based anomaly explanation, showing that retrieval-augmented models outperformed pure generative baselines on factual accuracy by a significant margin. In SYNAPSE, RAG is operationalized at two distinct stages of the pipeline: first, at the DECI causal discovery stage where historical pattern embeddings are retrieved to construct the $W_0$ initialization prior that biases causal graph learning toward structurally validated topologies; and second, at the LLM reasoning stage where relevant past incident summaries and service-specific runbook fragments are retrieved from a vector store to enrich the GPT-4o-mini prompt context. This dual application of RAG ensures that both the structural causal model and the natural-language explanations are continuously informed by the evolving operational history of the monitored system.\\\\

\hspace*{0.3in} A synthesis of the surveyed literature reveals that while substantial advances have been made in individual components of the AIOps stack — anomaly detection, causal discovery, LLM-based explanation, and continual learning — no existing open-source or commercial system integrates these capabilities into a cohesive, end-to-end automated pipeline. Commercial tools such as Dynatrace and Datadog offer production-grade observability but sacrifice extensibility and formal causal grounding. Academic GNN and causal inference approaches achieve high diagnostic precision in controlled benchmarks but lack the operational integration required for real-time production deployment. LLM-based reasoning work remains largely decoupled from structured causal evidence, increasing hallucination risk. Continual learning has been studied extensively in the ML literature but has seen negligible adoption in deployed AIOps systems. SYNAPSE addresses this multi-dimensional gap by architecting a unified ten-stage AI pipeline in which anomaly detection, severity classification, causal discovery, root cause extraction, LLM reasoning, and recovery action selection are composable, co-trained modules sharing a common data representation layer. The integration of EWC and MAML-based continual learning further distinguishes SYNAPSE from all surveyed work, enabling the platform to remain operationally accurate in dynamic, long-running production environments without disruptive model replacement cycles.\\\\
